
\begin{theorem}\label{fga-the-5-8}\dcref{5.8}\source{fga-explained:page-0130}\uses{}
\begin{verbatim}
([EGA] III 7.7.8,7.7.9)Let 5 be a noetherian scheme and
TT:X ?> 5 a projectivemorphism. Let S and T be coherentsheaveson X. Consider
the set-valuedcontravariantfunctorS)om{E,!F) on 5-schemes, which associatesto
any T -^ S the set of allO^T-li^^^^ homomorphisms HomxT{^T^^T)       where Et
and Tt denote the pull-backsof E and T under the projectionXt ?> X. IfJ^ is
flatover 5, then the above functorisrepresentableby a linearscheme V over S.
\end{verbatim}
\end{theorem}
\begin{proof}
\begin{verbatim}
Firstnote that if<^is a locallyfreeO^-niodule, then S)om{E,T) is
the functorT k^ H^{Xt, {T <S>Ox^^)t)- The sheafT ?o^ E^'is again flatover
5, so we can apply Theorem 5.7 to get a coherent sheaf Q, such that we have
5.3.SEMI-CONTINUITY AND BASE-CHANGE                    121

TT^{!F?o^ E^ (8H^ TT*^)= HorriQ^iQ.Q) forallquasi-coherentsheavesQ on S. In
          if/ :Spec R-^ S isany morphism then taking Q = J^Or we get
particular,
       Mor5(Spec/^,SpecSym(^^ Q) = Homo^ -rnod{Q,I^Or)

                                     = H\Xr,{T?o.E'^)r)
                                     = HomxR{ER,TR).
This shows that V = SpecSym^^ Q isthe requiredhnear scheme when E is
      freeon X. More generallyforan arbitrarycoherent E^ over any affineopen
locally
U C S there existvectorbundles Ei and Eq on Xjj and a rightexact sequence
El -^ Eq -^ E -^ 0. (This iswhere we need projectivity   of X ?> 5. Instead,we
could have assumed justpropernesstogetherwith the conditionthat locallyover S
we have such a resolutionof <^.)Then applyingthe above argument to the functors
S)om{Ei^T) and S)om{E^^T)^ we get coherentsheaves Qi and Qo on t/,and from
the naturaltransformationS)om{Eo,T) ?> S)om{Ei,!F) induced by the homomor-
phism El -^ Eq, we get a homomorphism Qi ?> Qo- Let Qu be itscokernel,and
put Yu = SpecSym^^ Qu. It followsfrom itsdefinition(and the leftexactness
of Hom) that the scheme Yu has the desireduniversalproperty over U. Therefore
allsuch Vt/, as U variesover an affineopen coverof 5, patch togetherto give the
desiredlinearscheme V. (In sheaf terms, the sheaves Qu willpatch togetherto
give a coherentsheaf Q on 5 with V = Spec Sym^^ Q.)                          D
\end{verbatim}
\end{proof}


\begin{remark}\label{fga-rem-5-9}\dcref{5.9}\source{fga-explained:page-0131}\uses{}
\begin{verbatim}
In particular,    note that the zero sectionVo C V iswhere the
universalhomomorphism vanishes. If / G Homxri^T^-^T) definesa morphism
(y9/: T ?> V, then the inverseimage /~^Vo isa closedsubscheme T' of T with
the universalproperty that ift/ ?> T is any morphism of schemes such that the
pull-backof / iszero,then U -^ T factorsvia T\
    5.3.4. Base-change for flat sheaves. The followingis the main resultof
Grothendieckon base change forflatfamiliesof sheaves,which isa consequence of
\end{verbatim}
\end{remark}
